% chapters/ch33_theoretical_foundation_v4.tex
% Chapter 33 (Part V): NRMO Theoretical Foundation v4
% Source: NRMO_v5_updated.pdf, Chapter 32 v4 paper (lines 10336-10593 of v5_layout.txt)

\chapter{NRMO Theoretical Foundation (v4)}
\label{ch:theoretical-foundation-v4}


\index{Theorem!Ruin Dilution}
\index{Theorem!CMDP Separation}
\index{CVaR}
\index{NRMO!Theoretical Foundation}
\nrmobox{Document identification}{
\textbf{Title}: NRMO Part~III: Theoretical Foundation\\[0.3em]
\textbf{Subtitle}: Non-Ruin Maximization Objective (NRMO)\\[0.3em]
\textbf{Author}: Takashi Ikeya\\[0.3em]
\textbf{Status}: Mathematically Corrected v4 | Integrated into
Unified Specification | 2026
}

\paragraph{Note on relationship to Part~VII.}
This chapter is the theoretical-foundation paper as integrated into
the v4 monograph. Part~VII (Mathematical Formalisation) is the
arXiv-format paper. Both contain the same four theorems (T1--T4)
but with slightly different framing: this chapter integrates them
into the monograph's narrative; Part~VII presents them in
publication form with full proofs.

% =====================================================================
\section{Abstract}
\label{sec:tf-v4-abstract}

We prove four results for NRMO:

\begin{itemize}
\item \textbf{(T1) Ruin Dilution} --- ruin probability $q$ bounds
unconditional time-average reward by $(1 - q) R_{\max}$.
\item \textbf{(T2) CMDP Separation} via explicit two-state MDP with
exact finite-horizon ruin formula.
\item \textbf{(T3) Exponential threshold family} achieves variational
supremum of $\Sigma(f)$ at rate $O(\Delta / \kappa)$.
\item \textbf{(T4) Polynomial-time CVaR tractability} for log-concave
distributions.
\end{itemize}

\noindent
Three errors from prior drafts are corrected and documented. All
theorem statements and proofs are self-consistent. Experiments
across three domains show 57--65\% ruin reduction versus approximate
baselines.

% =====================================================================
\section{Introduction}
\label{sec:tf-v4-introduction}

Peters (2019) shows that for multiplicative growth processes, any
policy with positive ruin probability converges to ruin with
probability one under repeated play. This motivates NRMO's core
design: \emph{ruin avoidance as a hard constraint, not a soft
penalty}.

Standard CMDPs (Altman 1999) guarantee only \emph{expected}
constraint satisfaction, permitting ruin in individual trajectories.
NRMO enforces a \emph{trajectory-level} ruin-free constraint.

\paragraph{Contributions.}

\begin{itemize}
\item \textbf{(C1)} Theorem~1: Ruin Dilution --- $U(\pi) \le (1 -
q) R_{\max}$.
\item \textbf{(C2)} Lemma~1: Constructive feasibility of
$\mathrm{safe}(T, \varepsilon)$.
\item \textbf{(C3)} Theorem~2: CMDP Separation via explicit
two-state MDP.
\item \textbf{(C4)} Theorem~3: Variational supremum of exponential
threshold family [T3 corrected in v4].
\item \textbf{(C5)} Theorem~4: Polynomial-time CVaR tractability.
\item \textbf{(C6)} Empirical evaluation with explicitly approximated
baselines.
\end{itemize}

% =====================================================================
\section{Related Work}
\label{sec:tf-v4-related}

CMDPs (Altman 1999) and Lagrangian methods (Tessler et al.~2018)
guarantee \emph{expected} constraint satisfaction. CPO (Achiam et
al.~2017) provides trust-region updates with constraint guarantees
at each iterate. These methods permit ruin in individual
trajectories; NRMO enforces \emph{trajectory-level} safety.

Risk-sensitive RL (Chow and Pavone 2014) minimises CVaR; Theorem~4
formalises the connection. Non-ergodic economics (Peters 2019;
Kelly 1956) motivates prioritising time-average over
ensemble-average. Robust MDPs (Iyengar 2005; Nilim and El Ghaoui
2005) handle uncertain kernels; NRMO's robust chance constraint uses
this structure.

% =====================================================================
\section{Problem Formulation}
\label{sec:tf-v4-formulation}

Let $M = (S, A, P, r, S_F)$ be an MDP:

\begin{itemize}
\item $S$ finite.
\item $r: S \times A \to [-R_{\max}, R_{\max}]$ bounded.
\item $S_F \subset S$ absorbing failure states.
\end{itemize}

\begin{definition}[Failure absorption]
\label{def:failure-absorption}
$s_F \in S_F$ satisfies $P(s_F \mid s_F, a) = 1$, $r(s_F, a) = 0$
for all $a$. First hitting time:
$\tau_F := \inf\{t \ge 0 : s_t \in S_F\}$.
\end{definition}

\begin{definition}[Safe policy class]
\label{def:safe-policy}
$\mathrm{safe}(T, \varepsilon) := \{\pi : P_\pi(\tau_F \le T) \le
\varepsilon\}$.
\end{definition}

\begin{definition}[Canonical objective, per Addendum U-DEF-01-A]
\label{def:tf-v4-objective}
\begin{equation}
U(\pi) := \liminf_{T \to \infty} \frac{1}{T}
\mathbb{E}_\pi\!\left[
\sum_{t=0}^{T-1} r(s_t, a_t)
\right] \qquad \text{(unconditional time-average).}
\end{equation}
\end{definition}

% =====================================================================
\section{Main Theoretical Results}
\label{sec:tf-v4-main-results}

\subsection{Theorem 1: Ruin Dilution}

\begin{theorem}[Ruin Dilution]
\label{thm:tf-v4-ruin-dilution}
Let $\pi$ satisfy $P_\pi(\tau_F < \infty) = q \in [0, 1)$. Then:
\begin{enumerate}
\item[(i)] $U(\pi) \le (1 - q) \cdot R_{\max}$.
\item[(ii)] Any $\pi'$ with $U(\pi') > (1 - q) R_{\max}$ strictly
dominates $\pi$.
\end{enumerate}
\end{theorem}

\begin{proof}
\medskip\noindent\textbf{Part (i).}
Decompose by $\tau_F < \infty$ (probability $q$) and
$\tau_F = \infty$ (probability $1 - q$):
\begin{equation}
\mathbb{E}_\pi\!\left[\frac{1}{T} \sum r_t\right]
=
q \cdot \mathbb{E}\!\left[
\frac{1}{T} \sum r_t \mid \tau_F < \infty
\right]
+ (1 - q) \cdot \mathbb{E}\!\left[
\frac{1}{T} \sum r_t \mid \tau_F = \infty
\right].
\end{equation}

\textit{Ruin term}: On $\tau_F < \infty$, $r_t = 0$ for all $t \ge
\tau_F$ (Definition~\ref{def:failure-absorption}), so
$\mathbb{E}[(1/T) \sum r_t \mid \tau_F < \infty] \le R_{\max} \cdot
\mathbb{E}[\tau_F / T \mid \tau_F < \infty] \to 0$ as $T \to \infty$
($\mathbb{E}[\tau_F \mid \tau_F < \infty] < \infty$ since $S$ is
finite).

\textit{Survival term}: $\mathbb{E}[(1/T) \sum r_t \mid \tau_F =
\infty] \le R_{\max}$.

Combining: $U(\pi) = \liminf \le \limsup \le q \cdot 0 + (1 - q)
\cdot R_{\max}$.

\medskip\noindent\textbf{Part (ii).}
Follows immediately from Part~(i): $U(\pi') > (1 - q) R_{\max} \ge
U(\pi)$.
\end{proof}

\begin{remark}[Scope]
Part~(ii) requires the existence of $\pi'$ with $U(\pi') > (1 - q)
R_{\max}$. When safe actions yield low reward, this may not hold.
The theorem quantifies the \emph{ceiling} on any positive-ruin
policy.
\end{remark}

\subsection{Lemma 1: Feasibility of safe}

\begin{lemma}[Feasibility of safe]
\label{lem:tf-v4-feasibility}
Assume:
\begin{itemize}
\item[(A1)] $S$ finite;
\item[(A2)] $\exists \pi^{\mathrm{hold}}$ with
$P_{\pi^{\mathrm{hold}}}(\tau_F \le T) = 0$ for all $T$;
\item[(A3)] $\varepsilon > 0$.
\end{itemize}
Then $\mathrm{safe}(T, \varepsilon) \ne \emptyset$.

Sufficient condition for (A2):
$\exists s^* \in S \setminus S_F, a^*$ with $P(s^* \mid s^*, a^*) =
1$.
\end{lemma}

\begin{proof}
Under (A2), $P_{\pi^{\mathrm{hold}}}(\tau_F \le T) = 0 \le
\varepsilon$, so $\pi^{\mathrm{hold}} \in \mathrm{safe}(T,
\varepsilon)$. Sufficient condition: $\pi^{\mathrm{hold}}(s) = a^*$
keeps all trajectories at $s^* \notin S_F$.
\end{proof}

\subsection{Theorem 2: CMDP Separation}

\paragraph{Correction note.}
Prior v2 construction ($s_0 \leftrightarrow s_1$ recurrent loop) was
flawed: for any $\delta > 0$, long-run ruin probability $= 1$
(standard absorbing chain result). Formula
$q(\delta) = (1 - p) \delta / (1 - p \delta)$ computed only
single-visit ruin probability, not eventual. Corrected: two-state
MDP with independent per-step ruin events.

\begin{theorem}[NRMO Strictly Separates from Lagrangian CMDP]
\label{thm:tf-v4-cmdp-separation}
There exists MDP $M^*$ and $b \in (0, 1)$ such that:
\begin{enumerate}
\item[(i)] the primal CMDP optimum $\pi^*$ satisfies
$P_{\pi^*}(\tau_F \le T) = b > 0$ (constraint binds at optimum by
KKT);
\item[(ii)] $\mathrm{safe}(T, 0) \ne \emptyset$, so NRMO achieves
zero ruin.
\end{enumerate}
\end{theorem}

\begin{proof}
\medskip\noindent\textbf{Construction $M^*$.}
$S = \{s_0, s_F\}$, $A = \{a_{\mathrm{safe}}, a_{\mathrm{risky}}\}$.
$R = 5$, $p = 0.7$, $T = 10$, $b = 0.20$. From $s_0$:
$a_{\mathrm{risky}} \to s_F$ w.p.~$1 - p = 0.3$, stay w.p.~$p = 0.7$,
$r = R$. From $s_0$: $a_{\mathrm{safe}} \to$ stay, $r = 0$. $s_F$
absorbing, $r = 0$.

\medskip\noindent\textbf{Finite-horizon ruin formula.}
Under mixed policy $\delta = P(a_{\mathrm{risky}} \mid s_0)$, each
step $s_0 \to s_F$ occurs with probability $(1 - p) \delta$
independently:
\begin{equation}
\Phi(\delta, T) = P(\tau_F \le T) = 1 - (1 - (1 - p) \delta)^T.
\end{equation}

\medskip\noindent\textbf{CMDP analysis.}
Unconstrained optimum is $\delta = 1$:
$\Phi(1, 10) = 1 - (0.7)^{10} = 0.972 > b = 0.20$. Constraint
violated. Since $\mathbb{E}[\mathrm{return}]$ and $\Phi$ are both
strictly increasing in $\delta$, complementary slackness forces
$\Phi(\delta^*, T) = b$ exactly at the optimum. Solving:
\begin{equation}
\delta^* = \frac{1 - (1 - b)^{1/T}}{1 - p}
= \frac{1 - (0.80)^{0.1}}{0.3} = 0.0744.
\end{equation}
Verify: $\Phi(0.0744, 10) = 1 - (0.9777)^{10} = 0.200 = b > 0$.
$\lambda^* = (\mathrm{d}\mathbb{E} / \mathrm{d}\delta) /
(\mathrm{d}\Phi / \mathrm{d}\delta) \rvert_{\delta = \delta^*} > 0$
is the unique KKT multiplier.

\medskip\noindent\textbf{NRMO.}
$\varepsilon = 0$ forces $\delta = 0$; $\Phi(0, T) = 0$.
Lemma~\ref{lem:tf-v4-feasibility} applies
($a_{\mathrm{safe}} = $ hold action).
\end{proof}

\begin{remark}
CMDP saturates the ruin budget $b$: it trades every unit of permitted
ruin for reward. NRMO forfeits risky reward to achieve zero ruin.
$U(\pi_{\mathrm{safe}}) = 0$ vs.~$U(\pi^*) = \delta^* R \cdot
(\text{survival factor}) > 0$. This trade-off is domain-dependent
and explicit.
\end{remark}

\subsection{Theorem 3: Variational Analysis (Corrected in v4)}

\paragraph{Correction note.}
Prior v2 proof applied Euler-Lagrange to
$\Sigma(f) = \int f'(x)(1 - x)\,\mathrm{d}x$, incorrectly yielding a
logarithmic interior maximiser. Prior v3 theorem statement said
``equivalent to minimising $\int f(x)\,\mathrm{d}x$'' --- wrong sign.
Correction: integration by parts gives $\Sigma(f) = -\varepsilon_{\mathrm{floor}}
+ \int f(x)\,\mathrm{d}x$, so maximising $\Sigma(f) \Leftrightarrow$
maximising $\int f(x)\,\mathrm{d}x$. No smooth interior maximiser
exists.

\begin{theorem}[Exponential Family Achieves Variational Supremum,
v4 corrected]
\label{thm:tf-v4-variational}
Let $\mathcal{F} = \{f : [0, 1] \to [\varepsilon_{\mathrm{floor}},
\varepsilon_{\mathrm{global}}]:$ monotone non-decreasing,
$f(0) = \varepsilon_{\mathrm{floor}}$,
$f(1) = \varepsilon_{\mathrm{global}}$, $f \in C^1\}$. Define
$\Sigma(f) = \int_0^1 f'(x)(1 - x)\,\mathrm{d}x$.
\begin{enumerate}
\item[(i)] By integration by parts:
$\Sigma(f) = -\varepsilon_{\mathrm{floor}} +
\int f(x)\,\mathrm{d}x$, so maximising $\Sigma(f)$ is equivalent to
maximising $\int f(x)\,\mathrm{d}x$ over $\mathcal{F}$.
\item[(ii)] $\sup_{f \in \mathcal{F}} \int f(x)\,\mathrm{d}x =
\varepsilon_{\mathrm{global}}$, approached but not achieved by
$f_\kappa(x) = \varepsilon_{\mathrm{floor}} + \Delta(1 - e^{-\kappa
x}) / (1 - e^{-\kappa})$ as $\kappa \to \infty$.
\item[(iii)] $\varepsilon_{\mathrm{global}} - \int f_\kappa =
O(\Delta / \kappa)$.
\end{enumerate}
\end{theorem}

\begin{proof}
\medskip\noindent\textbf{Part (i): Integration by parts.}
$\Sigma(f) = \int_0^1 f'(x)(1 - x)\,\mathrm{d}x$. Set $u = (1 - x)$,
$\mathrm{d}v = f'(x)\,\mathrm{d}x$:
\begin{align}
\Sigma(f) &= [(1 - x) f(x)]_0^1 + \int_0^1 f(x)\,\mathrm{d}x \\
&= (0 \cdot f(1) - 1 \cdot f(0)) + \int f(x)\,\mathrm{d}x \\
&= -\varepsilon_{\mathrm{floor}} + \int f(x)\,\mathrm{d}x.
\end{align}
Since $-\varepsilon_{\mathrm{floor}}$ is constant: $\max \Sigma(f)
\Leftrightarrow \max \int f(x)\,\mathrm{d}x$.

\medskip\noindent\textbf{Part (ii): No smooth interior maximiser.}
For monotone $f$ with $f(0) = \varepsilon_{\mathrm{floor}}$,
$f(1) = \varepsilon_{\mathrm{global}}$, maximising
$\int f(x)\,\mathrm{d}x$ means rising to
$\varepsilon_{\mathrm{global}}$ as fast as possible from $x = 0$.
The supremum is achieved by $f^*(x) = \varepsilon_{\mathrm{floor}}
\cdot \mathbf{1}_{x = 0} + \varepsilon_{\mathrm{global}} \cdot
\mathbf{1}_{x > 0}$ (step function), giving $\int f^* =
\varepsilon_{\mathrm{global}}$. But $f \in C^1$, so no smooth
maximiser exists. $f_\kappa$ approximates $f^*$: for any $x > 0$,
$f_\kappa(x) \to \varepsilon_{\mathrm{global}}$ as $\kappa \to
\infty$.

\medskip\noindent\textbf{Part (iii): Convergence rate.}
\begin{equation}
\int_0^1 f_\kappa(x)\,\mathrm{d}x =
\varepsilon_{\mathrm{floor}} + \frac{\Delta}{1 - e^{-\kappa}}
\cdot \!\left[1 - \frac{1 - e^{-\kappa}}{\kappa}\right]
= \varepsilon_{\mathrm{floor}} + \Delta\!\left[1 - \frac{1}{\kappa}
+ O(e^{-\kappa})\right].
\end{equation}
Therefore $\varepsilon_{\mathrm{global}} - \int f_\kappa = \Delta /
\kappa + O(\Delta e^{-\kappa}) = O(\Delta / \kappa)$.
\end{proof}

\begin{remark}
The exponential family is the canonical smooth approximation to the
(unattainable) step-function supremum. Larger $\kappa$ is always
variationally better; $\kappa = 3$ is a practical balance against
numerical instability near $x = 0$. For $\kappa = 3$:
$\int f_3 - \varepsilon_{\mathrm{floor}} \approx 0.667 \Delta$,
leaving 33\% sensitivity gain unrealised relative to the limit.
\end{remark}

\subsection{Theorem 4: CVaR Tractability}

\begin{theorem}[Polynomial-Time CVaR Bound]
\label{thm:tf-v4-cvar}
Let $Z = \mathbf{1}_{\tau_F \le T} \in \{0, 1\}$. For any $\alpha \in
(0, 1)$:
\begin{enumerate}
\item[(i)] $\mathrm{CVaR}_\alpha(Z) \ge P(Z = 1)$.
\item[(ii)] $\min_\pi \mathrm{CVaR}_\alpha(\mathrm{cost}_\pi)$
subject to $\mathrm{CVaR} \le \varepsilon$ is solvable in
$O(|S|^2 |A| T \cdot \log(1 / \varepsilon))$ when the loss
distribution is log-concave.
\end{enumerate}
\end{theorem}

\begin{proof}
\medskip\noindent\textbf{Part (i).}
$Z \in \{0, 1\}$:
\begin{itemize}
\item Case A ($\alpha \ge 1 - P(Z = 1)$): $\mathrm{VaR}_\alpha(Z) =
1$, $\mathrm{CVaR} = 1 \ge P(Z = 1)$.
\item Case B ($\alpha < 1 - P(Z = 1)$): $\mathrm{VaR}_\alpha(Z) =
0$, $\mathrm{CVaR}_\alpha(Z) = P(Z = 1) / (1 - \alpha) \ge
P(Z = 1)$.
\end{itemize}
Both cases: $\mathrm{CVaR} \ge P(Z = 1)$.

\medskip\noindent\textbf{Part (ii).}
Rockafellar-Uryasev (2000):
\begin{equation}
\mathrm{CVaR}_\alpha(Z) = \min_v\!\left\{v + \frac{1}{1 - \alpha}
\mathbb{E}[\max(Z - v, 0)]\right\}.
\end{equation}
In occupancy-measure LP (Altman 1999, Ch.~2),
$\mathbb{E}_\pi[\max(Z - v, 0)]$ is linear in occupancy measures
$\mu(s, a, t)$ for fixed $v$. LP has $|S| |A| T + 1$ variables;
interior-point methods: $O(|S|^2 |A| T \cdot \log(1 / \varepsilon))$.
The LP formulation is exact by construction (Altman 1999, Ch.~2);
log-concavity is used only to confirm the loss distribution is
unimodal.
\end{proof}

% =====================================================================
\section{Algorithm: Phased Parallel Execution}
\label{sec:tf-v4-algorithm}

Phased Parallel Execution resolves the contradiction between full
parallelism (Phase~1) and Type ZERO suppression authority (Phase~3
only). DEADLOCK handler (Step~8) outputs $\mathrm{SAFE\_EXIT} =
\pi^{\mathrm{hold}}$ from Lemma~\ref{lem:tf-v4-feasibility}.

\begin{lstlisting}[basicstyle=\ttfamily\small,frame=single]
Algorithm 1: NRMO Execution

Input:  M, T, eps_global, eps_floor, kappa, B_t
Output: a_t in A or SAFE_EXIT

1: eps_t = eps_floor + (eps_global - eps_floor)
            * (1 - exp(-kappa * B_t))
2: PHASE 1 [parallel, read-only]:
     each module m computes score_m(s_t)
3: PHASE 2 [barrier]:
     await; T_max cutoff; late scores set to 0
4: PHASE 3 [serial]:
     Priority Arbiter -> Tri-Option Generator -> {A, B, C}
5: For o in {A, B, C}:
     verify APSOC gates
       (reversibility, Max Loss, exit trigger)
6: Compute CVaR upper bound on P(tau_F <= T)
     per option (Theorem 4)
7: O_valid = {o : CVaR <= eps_t and APSOC passed}
8: if O_valid == empty:
     DEADLOCK -> return SAFE_EXIT (pi_hold)
9: else:
     return argmax_{o in O_valid} U(pi_o)
\end{lstlisting}

% =====================================================================
\section{Experiments}
\label{sec:tf-v4-experiments}

\paragraph{Correction note.}
CPO and RCPO are approximated as analytic greedy Lagrangian policies
($\lambda_{\mathrm{CPO}} = 2.0$, $\lambda_{\mathrm{RCPO}} = 4.0$).
Results are directional evidence only. Faithful library
implementations deferred to future work.

\begin{table}[ht]
\centering
\scriptsize
\begin{tabularx}{\linewidth}{lXXXXXXXX}
\toprule
\textbf{Method} & \textbf{D1 Rwd} & \textbf{D1 Ruin$\downarrow$}
& \textbf{D2 Ret} & \textbf{D2 Ruin$\downarrow$} & \textbf{D2 SR}
& \textbf{D3 Profit} & \textbf{D3 Ruin$\downarrow$}
& \textbf{D3 Svc} \\
\midrule
\textbf{NRMO} & \textbf{8.22} & \textbf{0.106} & $-9.0\%$
& \textbf{0.193} & $-0.49$ & \textbf{17{,}548} & 0.000 & 0.997 \\
CPO & 6.13 & 0.248 & $-11.3\%$ & 0.467 & $-0.60$ & 16{,}306 & 0.000
& 0.998 \\
RCPO & 8.06 & 0.118 & $-9.7\%$ & 0.250 & $-0.54$ & 17{,}176 & 0.000
& 0.995 \\
Lagrangian & 5.69 & 0.278 & $-13.8\%$ & 0.550 & $-0.77$
& 16{,}469 & 0.000 & 0.976 \\
\bottomrule
\end{tabularx}
\caption[Empirical results]{Results
($n = 500 / 300 / 300$). $\downarrow$ = lower better. NRMO row
highlighted. Baselines are analytic approximations.}
\label{tab:tf-v4-empirical}
\end{table}

\paragraph{D1 Grid World.}
NRMO ruin $0.106$ vs.~$0.278$ Lagrangian ($-62\%$),
$0.118$ RCPO ($-10\%$). Hard ruin filter rejects options above
$\varepsilon_t$ entirely.

\paragraph{D2 Portfolio.}
Buffer-adaptive allocation reduces ruin to $0.193$ vs.~$0.550$
Lagrangian ($-65\%$). As $B_t \to 0$, $\varepsilon_t \to
\varepsilon_{\mathrm{floor}}$ and risky allocation decreases.

\paragraph{D3 Inventory.}
All methods zero ruin; NRMO leads profit ($+7.6\%$ vs.~CPO) via
emergency restock trigger.

% =====================================================================
\section{Discussion}
\label{sec:tf-v4-discussion}

\paragraph{T1 scope.}
Does not assert NRMO achieves higher reward than CMDP in general.
When safe actions earn zero reward, $U(\pi_{\mathrm{safe}}) = 0$
while CMDP achieves $U > 0$.

\paragraph{T2 scope.}
Separation instance $M^*$ is minimal. Practically relevant when
$T \cdot \delta^* (1 - p) = O(1)$.

\paragraph{T3 scope.}
No unique optimal $\kappa$; larger always better variationally,
constrained by numerics.

\paragraph{T4 note.}
LP formulation is exact by construction (Altman 1999); log-concavity
confirms unimodality only.

\paragraph{Open questions.}

\begin{enumerate}
\item Continuous state spaces.
\item POMDP extension.
\item Online estimation.
\item Safety Gym comparison.
\end{enumerate}

% =====================================================================
\section{References}
\label{sec:tf-v4-references}

\begin{enumerate}
\item Achiam et al.~(2017). Constrained Policy Optimization.
\textit{ICML 2017}.
\item Altman, E.~(1999). \textit{Constrained Markov Decision
Processes}. Chapman \& Hall/CRC.
\item Chow \& Pavone (2014). Risk-sensitive MDPs.
\textit{IEEE CDC 2014}.
\item Chow et al.~(2018). Lyapunov-based safe RL. \textit{NeurIPS
2018}.
\item Garcia \& Fernandez (2015). Survey on safe RL.
\textit{JMLR} 16(1).
\item Iyengar (2005). Robust dynamic programming.
\textit{Math.~OR} 30(2).
\item Kelly (1956). Information rate. \textit{Bell Sys.~Tech.~J.}
35(4).
\item Nilim \& El Ghaoui (2005). Robust MDPs. \textit{OR} 53(5).
\item Peters (2019). Ergodicity problem. \textit{Nature Physics}
15(12).
\item Rockafellar \& Uryasev (2000). CVaR optimization.
\textit{J.~Risk} 2(3).
\item Tessler et al.~(2018). Reward constrained policy
optimization. \textit{ICLR 2019}.
\end{enumerate}

% =====================================================================
\section*{Connections to Other Parts}

\begin{itemize}
\item This chapter is the v4 corrected version of the theoretical
foundation. Part~VII presents the same content in arXiv-paper
form.
\item Theorem~\ref{thm:tf-v4-ruin-dilution} is the strongest theorem
linking $U$ and $q$; it is the technical backbone of the
``ruin avoidance is primary'' principle (Section~\ref{sec:invariants-unchanged}).
\item Algorithm~1 (Phased Parallel Execution) is the algorithmic
template for the Phased / Parallel OODA architecture
(Chapter~\ref{ch:parallel-ooda}).
\item The empirical validation in Section~\ref{sec:tf-v4-experiments}
is at the algorithmic level. Civilisation-scale empirical
validation is in Part~XI.
\end{itemize}
